% FANCY использует стиль plain для содержания и еще для чего-то
% здесь мы его переопределим, чтобы в Содержании отображались колонтитулы
\fancypagestyle{plain}{
\fancyhf{}
\renewcommand{\headrulewidth}{1pt}
\lhead{\color{red}{ЮТКБ.656122.611-01.01 РЭ}}
\chead{\color{black}{Редакция \rev~от~\revdate}}
\rhead{\color{unidarkgreen}\colondevicetype}
\lfoot{\color{unidarkgreen}{Руководство по эксплуатации }}
\rfoot{{\thepage}}
\renewcommand{\footrulewidth}{1pt}
}

\thispagestyle{plain} % подключаем переопределенный стиль plain к странице с содержанием

%\fancyhf{}
%\fancyhead{} % clear all header fields
%\renewcommand{\headrulewidth}{1pt}

%\rhead{ЮНИТ-М3-ДЗТ2}
%\lhead{ЮТКБ.656112.608 РЭ2}

%\renewcommand{\footrulewidth}{1pt}
%\lfoot{\hyperref[https://uni-eng.ru]{https://uni-eng.ru}}
%\rfoot {ООО «Юнител Инжиниринг»}

{\centering
\textcolor{unidarkgreen}{\Large\bfseries СОДЕРЖАНИЕ}\par
\vspace{-1.5em}
}
\renewcommand{\contentsname}{}
\tableofcontents

